%%%%
%%%% Initially nicked from:
%%%%  https://github.com/bow/talks/blob/master/2014_bosc/abstract/abstract.tex
%%%%
%%%% \url{http://geneontology.org/}
%%%%
%%%% Compile:
%%%%   pdflatex abstract.tex
%%%%   evince abstract.pdf
%%%%   TODO: pandoc -s -f latex -t markdown_github -o ../README.md ./abstract.tex
%%%%

\documentclass[10pt,oneside]{article}
\usepackage{enumitem} % remove spaces inside lists

%%%%%%%%%%%%%
\setlength{\textheight}{8.75in} %Letter is 11in, less 2 for margins, less 0.25 for footer
\setlength{\oddsidemargin}{0.0in} %gets +1inc
\setlength{\evensidemargin}{0.0in} %gets +1inch
\setlength{\textwidth}{6.50in} %Letter is 8.5, less 2 inches for margins
\setlength{\topmargin}{0.5in}
\setlength{\headheight}{0in}
\setlength{\headsep}{0in}
\setlength{\parindent}{0.25in}
%%%%%%%%%%%% 

\usepackage[numbers]{natbib}
\usepackage{graphicx}
\usepackage[colorlinks=true,citecolor=black,urlcolor=blue]{hyperref}

\title{%Hack to get the logo on the PDF front page:
\vspace{-1.5in}
%% Some latex cannot read PNG sizes.
\includegraphics[width=0.5\textwidth,natwidth=485,natheight=125]{go-logo-small.png}\\
%Hack to get some white space using a blank line:
~\\Users, Communication, and a Light Application-Level API: A Request for Comments}

\author{
  Seth Carbon\footnote{Berkeley Bioinformatics Open-source Projects, Lawrence Berkeley National Lab, Berkeley, CA, USA.} % Email: \href{mailto:sjcarbon@lbl.gov}{sjcarbon@lbl.gov}}
}

\date{17\textsuperscript{th} Bioinformatics Open Source Conference
  (BOSC) 2017, Prague, Czech Republic}

\begin{document}
%\pagestyle{empty}
%\thispagestyle{empty}
\maketitle
%\pagestyle{empty}
\thispagestyle{empty}

\vspace{-0.2in}
\noindent
Website: \url{https://github.com/kltm/lala} \\
Repository: \url{https://github.com/kltm/lala} \\
License: CC-BY 4.0 \\

%escaping many of the pitfalls of more ``tabular'' modeling. Noctua also
%The Gene Ontology \citep{GO} project aims to create a comprehensive

Much thought and paper has been put into sharing information between
resources, from the use of common identifiers to use of common data
stores and APIs. We would like to explore an area that has not, to the
authors' knowledge, been as well explored: the sharing of simple
information directly between web-based applications, without the use
of a non-application central server or API as an intermediary.
%% For example, a curator may wish to
%% import IDs from one tool to another, curate the same paper using
%% different tools, or a researcher may wish to analyze the same piece of
%% genome in different analysis tools.

More concretely, the use cases that we wish to look at are how to
manage user authentication, authorization, and how to obtain small
packets of specific information from an external resource that has its
own associated web application. As a way to explore this space, we
have created applications that accomplish this by passing tokens and
round-tripping a packet of information (acting as a ``black box'' or
``piggy bank'') through the external application using basic HTTP
methods, allowing an easy high-level ``federation''. In an environment
where much effort is being spent on the creation of rich web
applications tailored to the habits of specific groups of scientists
and curators, users and engineers would have much to gain by being
able to reuse functionality from external applications as-is in their
workflows and stacks, as well as having the added benefit of driving
traffic to external web applications that implement such a common
specification.

Through these explorations, we have decided that a path forward should
capture the following qualities:

%% For the prototype we worked on, we started with a method that Galaxy
%% [] uses for running a round-trip for interactively collecting data for
%% use in a pipeline. In addition to this pattern, we added an
%% information packet that is meant to be passed through the external
%% system with no modification. This model seemed like it would be
%% appropriate for our use case as it included the following properties:

\begin{itemize}[noitemsep]
\item Easy to implement: complexity can be a barrier to adoption and
  implementation
\item Basic HTTP tooling: any system should have easy access to
  the tools necessary
\item ``Stateless'': simplifying debugging and implementation
\item Minimal need for initial coordination \\ Beyond what data is to
  be returned, the external application does not need to understand
  the transiting packet
\item No need for calling application to coordinate changes to own API
  after initial coordination: as the calling application is
  responsible for decoding the information it initially sent and the
  location it is sent to, major API changes can occur without the need
  to coordinate with any other resource
\item Have the ability to perform operations ``remotely'' while still
  logged-in to the calling application, without the need to coordinate
  cross-site logins: this can be done by placing an authorization
  token in the transiting packet
\end{itemize}

Taking inspiration from the methods used by Galaxy \citep{Galaxy} to
pull data in from external applications, we'd like to discuss this
potential protocol and get feedback from the community about general
patterns for passing users and relatively simple data directly between
applications. Variations of this proposal have been implemented or
explored in applications such as: Noctua, Textpresso Central, AmiGO,
and PubAnnotation.

%By using a model that allows for
%deposit in, but no other alteration to, the data passed data packet,
%we are able to recreate state in the calling application to apply the
%external data to the desired task.

%% If the calling application and external application coordinate, or if
%% the calling application has a rich and documented pass-through packet
%% format, the external application can skim information from the
%% interaction, to create a richer interface for the user, or to add that
%% data to external application’s store.

%% This proposal has its origins in initial work done with Noctua [?]
%% communicating with Textpresso Central [?] and later work done with
%% looking at communication between Noctua and PubAnnotation [?] at
%% BioHackathon 16 \url{https://github.com/dbcls/bh16/wiki/Noctua#add-connections-with-external-markuppaper-annotation-resources-partial}
%%   ?] and BLAH 3 \url{https://docs.google.com/presentation/d/1-3HgoBmnu9p1Alt6myjfLUJIOojIgpskRFh47NjEd5A/edit?usp=sharing} \url{https://docs.google.com/document/d/16DNkAO9YHKZmQjdVxq2B_ULYQDu7jjc0aciVKdTwX8w/edit?usp=sharing}
%%   ?] to use external tools and IDs to enrich the evidence captured by
%% Noctua.

%% This project has two aspects: 1) the specific task of coordinating
%% basic communication between Noctua and PubAnnotation for adding PubMed
%% publication spans to evidence in Noctua models (and the possible
%% addition of Noctua modeling information to PubAnnotation held
%% annotations) and 2) to prototype a light model for sharing information
%% between two interactive applications--a calling application and an
%% external application.

% License 
%
% This work is a proposal for a software protocol, not a yet concrete or
% generalized implementation. As such, the text of this work is
% available under the CC-BY 4.0 license.

%## Current implementations and work

%% * https://github.com/geneontology/noctua/issues/283
%% * https://github.com/geneontology/noctua/issues/316
%% * https://github.com/pubannotation/pubannotation/issues/3

\begin{thebibliography}{}

%\bibitem[The Gene Ontology Consortium, 2015]{GO}Gene Ontology Consortium: going forward. {\it Nucl. Acids Res.} 43, D1049–D1056.
  
\bibitem[Enis Afgan, Dannon Baker, Marius van den Beek, Daniel Blankenberg, Dave Bouvier, Martin Čech, John Chilton, Dave Clements, Nate Coraor, Carl Eberhard, Björn Grüning, Aysam Guerler, Jennifer Hillman-Jackson, Greg Von Kuster, Eric Rasche, Nicola Soranzo, Nitesh Turaga, James Taylor, Anton Nekrutenko, and Jeremy Goecks, 2016]{Galaxy}The Galaxy platform for accessible, reproducible and collaborative biomedical analyses, 2016 update {\it Nucleic Acids Research} 44(W1): W3-W10, doi:10.1093/nar/gkw343
  
\end{thebibliography}

\end{document}
